In this section, we describe the design choices made to implement \scheme{} hardware.

\subsection{Changes in the Network Interface} \label{sec-ni-changes}
We extend the Network Interface (NI) to control the timing of packet ejection, which is central to the defense mechanism of \scheme{}.
NI extensions have previously been used to provide security features such as malware detection~\cite{novo2020flow} and network intrusion detection~\cite{mirsky2018kitsune}.
We extend NI capability to detect packets that require delay or bypass based on their expected latency, following the \scheme{} logic described in Section~\ref{sec-overview}.

When the NI receives a load request, it sets flags to decide between the two available networks (normal or bypass) or to add delay (for \jscheme{}).
The bypass network uses the same physical links as the normal network.
We use virtual channel (VC) separation to isolate bypass traffic from regular network traffic~\cite{duato2001general}.
Bypass packets use a dedicated VC, ensuring that bypass traffic does not interfere with normal network operation.

The NI uses a 1-bit flag \texttt{bypass\_flag} on each output port of the router.
When the NI sends any flit through an output port, it checks whether \texttt{bypass\_flag} is enabled for that cycle.
If the flag is disabled, the NI cannot send the flit on that cycle.
When \bscheme{} needs to send a flit via the bypass path, the NI first checks and sets this flag.
If the flag is already set, the NI stalls the flit and retries on subsequent cycles.

In the common case, the attacker targets only two distinct LLC banks.
Since XY routing produces non-overlapping bypass paths for different source-destination pairs in a mesh, bypass contention is unlikely in practice.
However, collisions can be resolved by pipelining the bypass mechanism.

When an NI sends a flit through either the bypass or regular path, it sets a timestamp in the header flit.
For request flits, the timestamp records the flit creation time.
For response flits, the timestamp records the creation time of the corresponding request flit.
This timestamp is used by the \scheme{} logic in the NI.
\scheme{} uses a comparator and multiplexer (MUX) in the NI to implement the defense.
Figure~\ref{fig:block-diagram} illustrates the bypass operation using the existing NI buffers.

\begin{figure}[t]
\centering
\includegraphics[width=\linewidth]{figs/NI Block Diagram 2.pdf}
\caption{Block diagram of bypass operations in the extended Network Interface.}
\label{fig:block-diagram}
\end{figure}

\subsection{Delay Implementation} \label{sub:jitter}
To implement dynamic delay, each NI includes a Delay Queue (DQ) that buffers flits associated with Security Sensitive instructions until their target latency is met.
Every flit carries a Target Latency (TL) field in its header, set to $T_{worst}$ (\jscheme{}), $T_{low}$ (\bscheme{}), or a pseudorandomly drawn $T_{rand}$ (\rscheme{}) based on the defense scheme in use.

For \rscheme{}, $T_{rand}$ is derived rather than stored.
Both the destination NI and the source NI compute $T_{rand}$ independently, using a pseudorandom function of the timestamp already carried in the flit header (Section~\ref{sec-ni-changes}).
Both NIs apply the same function with the same seed to the same timestamp, so they compute the identical value.
No extra communication or storage is needed beyond the existing timestamp field.

When the tail flit of a Security Sensitive response arrives at the source NI, the NI computes the Round-Trip Latency (RTL) of the request using the stored timestamp and compares it to TL using a comparator.
There are two cases:
\begin{compactenum}
    \item \textbf{RTL $\geq$ TL:} No delay is needed.
    The flit is forwarded to the ejection queue of the NI immediately, or placed in the stall queue if the ejection queue is unavailable.

    \item \textbf{RTL $<$ TL:} The flit is placed in the DQ at the source NI with a delay of $TL - RTL$ cycles.
    A counter wakes up the NI to eject the flit after the specified delay has elapsed.
\end{compactenum}
